% =====================================================================
%  MLE-AI-201 Machine Learning -- Final Project (Team 07)
%  Build:   pdflatex report.tex ; pdflatex report.tex
% =====================================================================
\documentclass[10pt,twocolumn,a4paper]{article}

\usepackage[a4paper,top=0.75in,bottom=1.0in,left=0.68in,right=0.68in,columnsep=0.25in]{geometry}
\usepackage{mathptmx}
\usepackage[T1]{fontenc}
\usepackage[utf8]{inputenc}
\usepackage{graphicx}
\usepackage{amsmath,amssymb}
\usepackage{booktabs}
\usepackage{array}
\usepackage{titlesec}
\usepackage{caption}
\usepackage{cite}
\usepackage[hidelinks]{hyperref}
\usepackage{url}

\renewcommand{\thesection}{\Roman{section}}
\renewcommand{\thesubsection}{\Alph{subsection}}
\titleformat{\section}[block]{\centering\normalsize\scshape}{\thesection.}{0.5em}{}
\titleformat{\subsection}[block]{\normalsize\itshape}{\thesubsection.}{0.5em}{}
\titlespacing*{\section}{0pt}{1.4ex plus 1ex minus .2ex}{1ex}
\titlespacing*{\subsection}{0pt}{1.0ex plus 1ex minus .2ex}{0.6ex}

\newcommand{\abstractkw}[1]{\noindent\textbf{\textit{#1}}}

\captionsetup{font=small,labelsep=period}
\setlength{\parindent}{1em}

\begin{document}

% ------------------------- TITLE BLOCK --------------------------------
\twocolumn[
  \begin{@twocolumnfalse}
  \begin{center}
    {\LARGE\bfseries Predicting Azerbaijani Property Prices from Scratch:\\
    A NumPy Decision Tree and Pegasos SVM on bina.az Listings\par}
    \vspace{1.2em}
    {\large
      Badalov, Sanan\quad Hasanli, Ali\quad Hasanli, Oghuz\quad Sadigov, Sadigov\par}
    \vspace{0.5em}
    {\itshape MLE-AI-201 Machine Learning, AI Academy --- Team 07\par}
    \vspace{1.4em}
  \end{center}
  \begin{quote}
    \small
    \abstractkw{Abstract---}%
    We turn the raw bina.az sale-listings dump (100{,}775 rows) into a leakage-free
    modelling set and frame it as two tasks: regression of $\log(\text{price})$
    (Task~A) and classification of a premium/standard price tier split at the
    training median (Task~B). We implement in NumPy a decision tree (Gini, entropy,
    MSE), a linear soft-margin SVM trained with Pegasos sub-gradient descent, and
    an RBF-like SVM using random Fourier features, and benchmark them against
    scikit-learn with matched settings. Hyper-parameters are chosen on a validation
    split only. On Task~A our tree reaches test $R^2=0.811$ (RMSE\,$=0.340$ in log
    space), identical to scikit-learn to four decimals. On Task~B our tree reaches
    F1\,$=0.930$, versus $0.793$ for the linear SVM and $0.850$ for the RFF-SVM;
    scikit-learn's kernel SVC reaches $0.868$. Tree-based models win because price
    depends on non-linear interactions of area and location. Code:
    \href{https://github.com/BadalovSanan/team07_ml_final_project}{\texttt{github.com/BadalovSanan/team07\_ml\_final\_project}}
    (tag \texttt{v1.0-final}).

    \vspace{0.6em}
    \abstractkw{Index Terms---}decision tree, support-vector machine,
    Pegasos, random Fourier features, regression, classification,
    real-estate price prediction.
  \end{quote}
  \vspace{1.0em}
  \end{@twocolumnfalse}
]

% =====================================================================
\section{Introduction}
Property prices are a standard tabular-learning problem, and the bina.az dump is
realistic: scraped, mixed-type, and full of outliers. Our goals were to
(i) clean it and remove target leakage, (ii) implement a decision tree and an SVM
\emph{from scratch}, and (iii) compare them honestly with scikit-learn~\cite{sklearn}.
Contributions: a documented cleaning pipeline; a tested NumPy tree supporting
Gini/entropy/MSE and four stopping rules; a Pegasos SVM~\cite{pegasos} with margin
and support-vector reporting plus a random-Fourier-feature extension~\cite{rff};
a from-scratch Ridge baseline; and a leakage-free benchmark with
cross-validated uncertainty, timing, and error analysis.

\section{Background}
A tree greedily picks the split maximising impurity decrease~\cite{cart}:
$G=1-\sum_k p_k^2$, $H=-\sum_k p_k\log_2 p_k$, and for regression
$I_\text{reg}=\frac1N\sum_i(y_i-\bar y)^2$, with gain
$IG=I(\text{parent})-\sum_c\frac{N_c}{N}I(c)$.
The soft-margin SVM minimises the regularised hinge loss, which Pegasos
solves by stochastic sub-gradient steps with rate $\eta_t=1/(\lambda t)$~\cite{pegasos}.

\section{Data and Preprocessing}
\textbf{Data.} The Kaggle ``binaaz-sale-project'' dump (downloaded 2026-10-14,
version~1) has 100{,}775 rows and 51 columns with Azerbaijani names and text.

\textbf{Cleaning} (all in \texttt{src/data\_prep.py}). Numeric columns arrive as
strings, so we parse area (``145 m$^2$''; ``sot'' converted at 1~sot\,$=$\,100~m$^2$),
land area, rooms, and floor (``7\,/\,9'' $\to$ floor, number of floors). All prices
are in AZN, so no conversion is needed. We drop rows with missing price, area, or
coordinates; remove content duplicates (same price, area, rooms, floor, coordinates,
category, location, and description, since the dump has no exact duplicate ids);
and apply \emph{fixed} outlier bounds (not data-derived, hence leak-free):
price in $[3\text{k},10\text{M}]$ AZN, area in $[10,50\text{k}]$~m$^2$, rooms\,$\le$\,20,
and coordinates inside Azerbaijan. Structurally missing floor, rooms, and land
values (a house has no floor) are filled with 0 plus a missing-indicator, so no
statistic is learned from the data. Azerbaijani columns we keep are renamed
(\texttt{category}, \texttt{repair}, \texttt{bill\_of\_sale}, \texttt{mortgage}).
Categoricals are one-hot encoded over a fixed category list, so every split has
identical columns.

\textbf{Leakage.} We drop \texttt{unit\_price} (price per m$^2$) and
\texttt{total\_price} (equal to \texttt{price} in 99.99\% of rows), both arithmetic
functions of the target. We also drop identifiers and scrape bookkeeping (ids,
URLs, owner and shop names, raw address, timestamps), post-listing popularity
and promotion fields (views, vip, featured), free text, currency (constant), and
a building-type column with only 154 non-null values.

\textbf{Tasks and split.} Task~A predicts $\log(\text{price})$ (prices are heavy
tailed, Fig.~\ref{fig:eda}); metrics are RMSE, MAE, and $R^2$, reported in log
space and in AZN. Task~B labels a listing \emph{premium} if its price exceeds the
\textbf{training-set median} (212{,}750~AZN), so the classes are balanced by
construction; metrics are precision, recall, F1, ROC-AUC, and the confusion matrix.
The 70/15/15 split is stratified on above/below median price, and standardisation
statistics come from the training split only (Table~\ref{tab:data}).

\begin{table}[t]
  \centering
  \caption{Dataset summary after cleaning.}
  \label{tab:data}
  \begin{tabular}{@{}lr@{}}
    \toprule
    Property & Value \\
    \midrule
    Rows raw $\to$ cleaned      & 100{,}775 $\to$ 75{,}660 \\
    Features used               & 21 \\
    Train / Val / Test split    & 52{,}964 / 11{,}348 / 11{,}348 \\
    Tier threshold (train median) & 212{,}750 AZN \\
    Price-tier balance (std/prem) & 50 / 50 \\
    \bottomrule
  \end{tabular}
\end{table}

\begin{figure}[t]
  \centering
  \includegraphics[width=\linewidth]{figures/eda_price_distribution.png}
  \caption{Price (left, clipped at the 99th percentile) is right-skewed;
  $\log(\text{price})$ (right) is close to symmetric, motivating Task~A on the log scale.}
  \label{fig:eda}
\end{figure}

\section{Method}
\subsection{Decision tree (from scratch)}
\texttt{DecisionTree} mirrors the scikit-learn API (\texttt{fit}, \texttt{predict},
\texttt{predict\_proba}). For each feature, values are sorted once and all midpoints
between distinct neighbours are scored in one vectorised pass using cumulative class
counts (classification) or cumulative sums and sums of squares (regression), so
the split search is $O(n\log n)$ per feature. The best (feature, threshold) by
impurity decrease is taken and the tree recurses. Leaves predict the majority class
(with \texttt{predict\_proba} from leaf frequencies) or the mean. Supported
hyper-parameters are \texttt{criterion} (Gini, entropy, MSE), \texttt{max\_depth},
\texttt{min\_samples\_split}, \texttt{min\_samples\_leaf}, and
\texttt{min\_impurity\_decrease}; unit tests check that each one changes the tree and
that splits match scikit-learn on tiny determined examples and on larger data
(\texttt{tests/test\_decision\_tree.py}). Impurity-decrease feature importances
are accumulated over all splits.

\subsection{Support-vector machine (from scratch)}
We minimise
\begin{equation}
  J(\mathbf{w},b)=\tfrac{\lambda}{2}\lVert\mathbf{w}\rVert^{2}
  +\tfrac{1}{n}\sum_{i=1}^{n}\max\!\big(0,\,1-y_i(\mathbf{w}^{\top}\mathbf{x}_i+b)\big)
\end{equation}
with labels $y_i\in\{-1,+1\}$. At step $t$ we draw an example; with
$\eta_t=1/(\lambda(t+t_0))$, the update is
$\mathbf{w}\leftarrow(1-\eta_t\lambda)\mathbf{w}+\eta_t y_i\mathbf{x}_i\,[y_i(\mathbf{w}^\top\mathbf{x}_i+b)<1]$,
and $b\leftarrow b+\eta_t y_i$ on margin violations (unregularised).
We use $t_0=1/\lambda$ so the first steps are $\approx1$ rather than $1/\lambda$,
which would make the weights explode for small $\lambda$ (textbook Pegasos is
$t_0=0$). Constant and $1/\sqrt{t}$ schedules, mini-batches, and epochs are also
exposed. Features are standardised first. After training we report
$\lVert\mathbf{w}\rVert$, margin width $2/\lVert\mathbf{w}\rVert$, and the support
vectors ($y_i(\mathbf{w}^\top\mathbf{x}_i+b)\le1$).

\textbf{Non-linearity.} We chose random Fourier features
(option~ii) because they reuse the same linear solver and cost $O(n)$ instead of an
$O(n^2)$ kernel matrix: $z(\mathbf{x})=\sqrt{2/D}\cos(W^\top\mathbf{x}+\mathbf{c})$,
$W\!\sim\!\mathcal N(0,2\gamma)$, $c\!\sim\!U(0,2\pi)$, $D=300$, so that
$z(\mathbf{x})^\top z(\mathbf{x}')\approx e^{-\gamma\lVert\mathbf{x}-\mathbf{x}'\rVert^2}$.
We also wrote a Ridge regression via the normal equations as a Task~A baseline.

\section{Experimental Setup}
\textbf{Protocol.} All seeds are fixed (\texttt{SEED=42}) and
\texttt{python -m src.run\_all} regenerates every table and figure. Hyper-parameters
are chosen on the \emph{validation} split only: tree depth
$\in\{1,\dots,20\}$ and \texttt{min\_samples\_leaf}$\in\{1,\dots,250\}$ (and
Gini vs.\ entropy for Task~B); SVM $\lambda\in\{10^{-6},\dots,10^{-1}\}$ and
$\gamma\in\{0.048,0.095,0.19,0.38\}$ for RFF; Ridge $\lambda\in\{10^{-3},\dots,10^{3}\}$.
The SVM uses 20 epochs, batch size~1. The test set is used once for final
numbers. Uncertainty is the mean\,$\pm$\,std of 5-fold cross-validation
(stratified for Task~B) on train\,$+$\,validation with the chosen settings.
Metrics are implemented in NumPy (\texttt{src/evaluate.py}).

\textbf{Baselines (matched).} Trees: \texttt{DecisionTreeRegressor/Classifier} with the
same criterion, depth, and leaf size. Linear SVM: \texttt{SGDClassifier}
(hinge, L2, \texttt{alpha}\,$=\lambda$, same epochs) in a \texttt{Pipeline} with
\texttt{StandardScaler}. Kernel SVM: \texttt{SVC} (RBF, same $\gamma$,
$C=1/(\lambda n)$) in a pipeline; because SVC is $O(n^2)$--$O(n^3)$, it is trained
on a random subset of at most 10{,}000 rows. We also include a mean (Task~A) and a
majority-class (Task~B) dummy baseline, and \texttt{sklearn.linear\_model.Ridge}.

\section{Results}
\textbf{Chosen configurations (validation only).} Task~A tree: MSE, depth 12,
\texttt{min\_samples\_leaf}\,$=5$. Task~B tree: Gini, depth 20,
\texttt{min\_samples\_leaf}\,$=1$ (Gini and entropy are within $0.001$ val F1:
$0.9328$ vs.\ $0.9324$). Linear SVM: $\lambda=10^{-3}$. RFF-SVM:
$\lambda=10^{-5}$, $\gamma=0.095$. Ridge: $\lambda=10^{-3}$.

\begin{table*}[t]
  \centering
  \small
  \caption{Task~A (regression of $\log$ price). Test-set metrics; CV is 5-fold mean\,$\pm$\,std on train+val; Fit in seconds on the full training split. AZN columns use $e^{\hat y}$.}
  \label{tab:resultsA}
  \begin{tabular}{@{}lcccccc@{}}
    \toprule
    Model & RMSE (log) & $R^2$ (log) & MAE (AZN) & RMSE (AZN) & CV RMSE (log) & Fit (s) \\
    \midrule
    Mean baseline              & 0.7806 & 0.000 & 180{,}660 & 504{,}971 & 0.7789 $\pm$ 0.0045 & 3.11 \\
    \textbf{Ours: tree (MSE)}  & \textbf{0.3397} & \textbf{0.811} & 78{,}773 & 276{,}905 & 0.3509 $\pm$ 0.0069 & 1.36 \\
    sklearn DecisionTreeRegressor & 0.3397 & 0.811 & 78{,}694 & 276{,}725 & 0.3505 $\pm$ 0.0065 & 0.40 \\
    Ours: Ridge ($\lambda{=}10^{-3}$) & 0.5801 & 0.448 & 163{,}061 & 911{,}242 & 0.5849 $\pm$ 0.0064 & 0.008 \\
    sklearn Ridge              & 0.5801 & 0.448 & 163{,}061 & 911{,}242 & 0.5849 $\pm$ 0.0064 & 0.26 \\
    \bottomrule
  \end{tabular}
\end{table*}

\begin{table*}[t]
  \centering
  \small
  \caption{Task~B (price tier, premium\,$=$\,positive). Test-set metrics; CV F1 is 5-fold mean\,$\pm$\,std on train+val; Fit/Pred in seconds (Pred on the test set).}
  \label{tab:resultsB}
  \begin{tabular}{@{}lcccccrr@{}}
    \toprule
    Model & Precision & Recall & F1 & ROC-AUC & CV F1 & Fit (s) & Pred (s) \\
    \midrule
    Majority-class baseline      & 0.000 & 0.000 & 0.000 & 0.500 & 0.000 $\pm$ 0.000 & 0.001 & 0.0001 \\
    \textbf{Ours: tree (Gini)}   & 0.928 & 0.932 & \textbf{0.930} & 0.932 & 0.931 $\pm$ 0.003 & 4.82 & 0.077 \\
    sklearn DecisionTreeClassifier & 0.930 & 0.933 & 0.931 & 0.934 & 0.931 $\pm$ 0.003 & 0.56 & 0.004 \\
    Ours: Pegasos SVM (linear)   & 0.782 & 0.804 & 0.793 & 0.868 & 0.790 $\pm$ 0.004 & 21.04 & 0.0003 \\
    sklearn SGDClassifier (hinge) & 0.782 & 0.807 & 0.794 & 0.868 & 0.790 $\pm$ 0.004 & 0.36 & 0.002 \\
    Ours: RFF-SVM (RBF approx.)  & 0.842 & 0.859 & 0.850 & 0.930 & 0.842 $\pm$ 0.005 & 20.96 & 0.078 \\
    sklearn SVC (RBF, $\le$10k rows) & 0.864 & 0.872 & 0.868 & 0.949 & 0.863 $\pm$ 0.003 & 3.29 & 5.00 \\
    \bottomrule
  \end{tabular}
\end{table*}

\textbf{Benchmark.} Our tree matches scikit-learn: identical test RMSE (log) and
$R^2$ on Task~A, and F1 within $0.0014$ on Task~B (Tables~\ref{tab:resultsA}
and~\ref{tab:resultsB}). The residual differences (MAE in AZN differs by $<0.2\%$)
come from tie-breaking between equal-gain splits and float rounding of midpoint
thresholds, not from different algorithms; the unit tests confirm identical splits
on determined examples. Our linear SVM matches \texttt{SGDClassifier}
(F1 $0.793$ vs.\ $0.794$, AUC $0.868$ for both), and our Ridge matches
scikit-learn's to all printed digits. Our implementations are slower to fit
(tree $3.4\times$ and $8.7\times$ on Tasks~A/B; SVM $58\times$) because the SVM
update is a per-example Python loop and scikit-learn uses compiled code. Kernel
SVC predicts slowly ($5.0$\,s vs.\ $0.08$\,s for RFF) because prediction scales with
the number of support vectors, whereas RFF prediction is a fixed-size map plus a dot product.

\begin{table}[t]
  \centering
  \small
  \caption{Linear Pegasos sweep over $\lambda$ (validation). Margin $=2/\lVert w\rVert$; SVs out of 52{,}964.}
  \label{tab:svm}
  \begin{tabular}{@{}rrrrrr@{}}
    \toprule
    $\lambda$ & $\lVert w\rVert$ & Margin & \#SV & Final $J$ & Val F1 \\
    \midrule
    $10^{-6}$ & 32.18 & 0.062 & 16{,}727 & 1.794 & 0.731 \\
    $10^{-4}$ & 4.35  & 0.460 & 28{,}604 & 0.520 & 0.785 \\
    $10^{-3}$ & 4.44  & 0.451 & 27{,}452 & 0.497 & \textbf{0.786} \\
    $10^{-2}$ & 2.53  & 0.790 & 30{,}830 & 0.554 & 0.786 \\
    $10^{-1}$ & 1.08  & 1.853 & 37{,}711 & 0.654 & 0.752 \\
    \bottomrule
  \end{tabular}
\end{table}

\textbf{Tree diagnostics.} Fig.~\ref{fig:depth} shows the bias--variance story.
Training error keeps falling with depth while validation error plateaus:
on Task~B, depth 20 gives train F1 $0.996$ but val F1 $0.933$ (depth 10:
$0.928/0.914$); on Task~A, depth 20 gives train RMSE $0.085$ but val
$0.371$, whereas depth 12 reaches val $0.364$. Raising \texttt{min\_samples\_leaf}
to 5 lowers Task~A val RMSE further to $0.353$. Learning curves (training-set
size 2.6k$\to$53k) show val F1 rising from $0.869$ to $0.933$ and val RMSE
falling from $0.458$ to $0.353$, still improving, so more data would help.
Impurity-decrease importances (Fig.~\ref{fig:imp}) put roughly half the weight on
\texttt{area\_m2} for both tasks, followed by \texttt{lat}/\texttt{lng}.

\begin{figure*}[t]
  \centering
  \includegraphics[width=0.49\linewidth]{{figures/WhatsApp Image 2026-10-18 at 00.55.58 (2).jpeg}}\hfill
  \includegraphics[width=0.49\linewidth]{{figures/WhatsApp Image 2026-10-18 at 00.55.58 (3).jpeg}}
  \caption{Tree validation curves. Left: F1 (Task~B) and RMSE (Task~A) vs.\ \texttt{max\_depth}; train and validation diverge beyond depth $\approx$10. Right: learning curves vs.\ training-set size.}
  \label{fig:depth}
\end{figure*}

\textbf{SVM diagnostics.} Table~\ref{tab:svm} shows the trade-off: a larger $\lambda$ widens the
margin ($0.062\to1.853$) and shrinks $\lVert w\rVert$, while val F1 peaks for
$\lambda\in[10^{-4},10^{-2}]$ and falls at both extremes (too small: not converged
in 20 epochs, $J=1.79$; too large: underfit). At $\lambda=10^{-3}$, 27{,}452 of
52{,}964 points ($52\%$) are support vectors, i.e.\ the classes are far from
linearly separable. The objective decreases over epochs (Fig.~\ref{fig:conv}):
$0.607\to0.497$ for the linear model and $0.516\to0.378$ for RFF, with noise from
stochastic updates (RFF spikes at epoch 10, then recovers). RFF improves F1 from
$0.793$ to $0.850$, showing non-linearity matters, and is sensitive to its
hyper-parameters (val F1 collapses to $0.53$ at $\gamma{=}0.38,\lambda{=}0.1$).
The linear weights (Fig.~\ref{fig:imp}) push towards premium for
\texttt{rooms}, \texttt{rooms\_missing}, and \texttt{area\_m2}, and towards
standard for land (\texttt{cat\_Torpaq}), \texttt{floor\_missing}, and house
categories. \texttt{rooms} and \texttt{rooms\_missing} are strongly
collinear (missing rooms are filled with 0), so their magnitudes should not be
read individually.

\begin{figure}[t]
  \centering
  \includegraphics[width=\linewidth]{{figures/WhatsApp Image 2026-10-18 at 00.55.58 (11).jpeg}}
  \caption{Pegasos objective $J$ per epoch for the linear SVM ($\lambda=10^{-3}$) and RFF-SVM ($\lambda=10^{-5}$).}
  \label{fig:conv}
\end{figure}

\begin{figure*}[t]
  \centering
  \includegraphics[width=0.49\linewidth]{{figures/WhatsApp Image 2026-10-18 at 00.55.58 (10).jpeg}}\hfill
  \includegraphics[width=0.40\linewidth]{{figures/WhatsApp Image 2026-10-18 at 00.55.58 (1).jpeg}}
  \caption{Left: impurity-decrease feature importances of our trees (Task~B, Task~A). Right: largest weights of the linear Pegasos SVM on standardised features.}
  \label{fig:imp}
\end{figure*}

\textbf{Error analysis.} Fig.~\ref{fig:err} breaks down errors. For Task~A the mean
absolute log-residual is lowest for apartments (\texttt{Yeni tikili} $0.138$,
\texttt{Kohne tikili} $0.135$) and highest for land plots (\texttt{Torpaq} $0.662$,
$n{=}544$) and commercial objects (\texttt{Obyekt} $0.494$, $n{=}485$); houses sit in
between ($0.268$). The worst test predictions are almost all land plots with
areas of $600$--$36{,}000$~m$^2$, where price depends on location and use that our
features do not capture (e.g.\ a 14{,}000~m$^2$ plot sold for 2.1M~AZN was predicted at
49k; a 36{,}200~m$^2$ plot at 230k was predicted at 5.0M). For Task~B the tree's
confusion matrix is $5262/412/384/5290$ (TN/FP/FN/TP). Error rate is $13.4\%$ for land
and objects but $5.9\%$--$6.1\%$ for apartments, and errors concentrate near the
threshold: $19\%$ in the closest quintile to 212{,}750~AZN versus $2$--$7\%$ elsewhere.
So the tree mostly fails where classes are genuinely ambiguous and in sparse,
heterogeneous categories (\texttt{Ofis} $n{=}17$, \texttt{Qaraj} $n{=}8$ are too small
to judge).

\begin{figure}[t]
  \centering
  \includegraphics[width=\linewidth]{{figures/WhatsApp Image 2026-10-18 at 00.55.58 (8).jpeg}}\\[2pt]
  \includegraphics[width=\linewidth]{{figures/WhatsApp Image 2026-10-18 at 00.55.58 (9).jpeg}}
  \caption{Error analysis of our tree. Top (Task~A): predicted vs.\ actual $\log$ price, residuals, and error by category. Bottom (Task~B): confusion matrix, error rate by distance to the tier threshold, and error rate by category.}
  \label{fig:err}
\end{figure}

\section{Discussion and Limitations}
\textbf{Why trees beat SVMs here.} Price is governed by interactions between
area, location, and property type, which axis-aligned threshold splits on
\texttt{lat}/\texttt{lng}/\texttt{area} capture directly; a linear boundary cannot
(F1 $0.793$), and Ridge on raw coordinates reaches only $R^2=0.448$ (its AZN-scale
$R^2$ is negative because exponentiating log predictions amplifies errors on
expensive listings). The RFF-SVM closes part of the gap ($0.850$) but trails
both our tree and scikit-learn's SVC ($0.868$, trained on only 10k rows). This suggests the
shortfall comes from the $D=300$ approximation and 20 epochs of stochastic
optimisation rather than from lack of data.

\textbf{Limitations.} (i) Results use one random split; CV runs on
train+val after tuning on val, so CV scores are slightly optimistic, while test numbers
are a single hold-out. (ii) The Task~B optimum sits at the edge of the depth grid
(20), so a deeper grid might help marginally. (iii) SVC uses at most 10k rows for
cost reasons, so it is not a like-for-like data budget. (iv) Our SVM is fixed at
20 epochs and is slow because of the per-sample Python loop. (v) Prices are
\emph{asking} prices from scraped listings; land, houses, and apartments share one
model, and location enters only through coordinates, with no district or time
features. (vi) We report ROC-AUC but not PR-AUC.

\section{Conclusion}
A from-scratch NumPy tree reproduces scikit-learn's accuracy on both tasks
($R^2=0.811$; F1\,$=0.930$), and a from-scratch Pegasos SVM reproduces
\texttt{SGDClassifier} with an inspectable margin and support-vector set. Trees
clearly outperform linear and RFF-SVMs on this data, and nearly all
errors come from land and commercial listings and from prices near the tier
threshold. Next steps would be district-level features, ensembles of our trees, and more
RFF components or kernelised Pegasos.

\section*{Tools and Acknowledgements}
An AI assistant (Claude, Anthropic) was used to draft the text of this report from
the repository's result files (\texttt{report/*.csv}, \texttt{*.md}) and figures;
all reported numbers come from those files and from \texttt{python -m src.run\_all}.
We built on NumPy, pandas, matplotlib, and scikit-learn (baselines only) and on the public
Kaggle dataset ``binaaz-sale-project'' by sehriyarmemmedli~\cite{data}.
Per-member contributions are listed in \texttt{contribution\_report.md}.

% =====================================================================
\begin{thebibliography}{9}
\bibitem{pegasos}
S.~Shalev-Shwartz, Y.~Singer, N.~Srebro, and A.~Cotter,
``Pegasos: primal estimated sub-gradient solver for SVM,''
\emph{Mathematical Programming}, vol.~127, no.~1, pp.~3--30, 2011.

\bibitem{rff}
A.~Rahimi and B.~Recht, ``Random features for large-scale kernel machines,''
in \emph{Advances in Neural Information Processing Systems (NeurIPS)}, 2007.

\bibitem{cart}
L.~Breiman, J.~Friedman, R.~Olshen, and C.~Stone,
\emph{Classification and Regression Trees}. Wadsworth, 1984.

\bibitem{sklearn}
F.~Pedregosa \emph{et al.}, ``Scikit-learn: machine learning in Python,''
\emph{Journal of Machine Learning Research}, vol.~12, pp.~2825--2830, 2011.

\bibitem{data}
S.~Memmedli, ``Bina.az Sale Project,'' Kaggle dataset, 2026.
[Online]. Available: \url{https://www.kaggle.com/datasets/sehriyarmemmedli/binaaz-sale-project}
\end{thebibliography}

\end{document}